\printedpage{43}
% head: A. Grothendieck : Dérivateurs, Ch. XII, Édition des Universités de Montpellier II et Paris VII ; foot: Version 7 décembre 2002

\msp{page 59}

\section*{V\quad Amalgamated sums and $W$-cocartesian squares in $\mathrm{Cat}$}

\subsection*{1\quad Generalities on $U \overset{i}{\hookrightarrow} X \overset{j}{\hookleftarrow} T$, and amalgamated sums.}

\subsubsection*{1.1\quad Gluing data between an open and a closed.}

Let $X$ be in Cat, $U$ an open in $X$, $T$ the closed complement. (\textbf{NB} The opens do indeed correspond to the opens of $\mathrm{Ob}\,X$, topologized by its natural preorder relation $x \leq y \iff \mathrm{Hom}_X(x,y) \neq \emptyset$, likewise for the closed ones. The topologies associated to preorder relations are those where every intersection of opens is still open, i.e. every union of closed sets is still closed. Thus the opens of such a topology are the closed sets of a dual topology, corresponding to the opposite preorder relation. The preorder relation associated to the topology is
\[
x \leq y \qquad\qquad \iff \qquad\qquad x \in \overline{\{y\}},
\]
the topology associated to the preorder on $I$ is the one for which the opens are the unions of sets of the form $I/x$ (which form a basis of opens), or again the \textit{``sieves''} ($x \in U$, $y \leq x \Longrightarrow y \in U$), the closed sets are the unions of sets $x\backslash I$, that is to say the cosieves. For every subset $A$ of $I$, there is a smallest open containing $A$, which is also its smallest neighborhood, it is the union of the $I/x$ for $x \in A$. And dually there is a largest closed set contained in $A$, whose complement is the smallest open containing $B = I\backslash A$, hence $\bigcup_{x\in B} I/x$. The closure of $A$ is the union of the closures $x\backslash I$ of its points.)

\msp{page 60}

The situation $(X,U,T)$ is reconstructed, up to unique isomorphism, from the knowledge of

\begin{enumerate}[label={}]
\item[$1^\circ$)] The two categories $U$, $T$.
\item[$2^\circ$)] The bifunctor $\displaystyle\left\{\begin{aligned}
&H : U^\circ \times T \longrightarrow \mathrm{Ens}\\[1ex]
&H(u,t) = \mathrm{Hom}_X(u,t).
\end{aligned}\right.$
\end{enumerate}

Moreover, these data $1^\circ$) $2^\circ$) are subject to no restriction. We have, more precisely % [sic] en sont -> ne sont, negation required by the sense

\textbf{Proposition 1.} \textit{An isomorphism} \texttt{[rather an equivalence of categories]} \textit{between the category $\mathcal{T}$ whose objects are the triples $(X,U,T)$, with $X \in \mathrm{Ob}\,\mathrm{Cat}$, $U \in \mathrm{Ouv}(X)$, $T \in \mathrm{Ferm}(X)$, $U$ and $T$ complementary to one another, the morphisms from $(X,U,T)$ to $(X',U',T')$ being the functors $f : X \longrightarrow X'$ which induce $U \longrightarrow U'$, $T \longrightarrow T'$ (whence
\printedpage{44}
% head: A. Grothendieck : Dérivateurs, Ch. XII, Édition des Universités de Montpellier II et Paris VII ; foot: Version 7 décembre 2002
$U = f^{-1}(U')$, $T = f^{-1}(T')$), with the category $\mathcal{T}'$ of triples $(U,T,H)$, of two categories $U$, $T$ and of a bifunctor}
\[
H : U^\circ \times T \longrightarrow \mathrm{Ens}
\]
\hspace*{2em}\textit{(or again a presheaf $U^\circ \longrightarrow T^\vee$ on $U$, hence an object of} $\uHom(U^\circ,T^\vee)$\textit{)),} % [sic] préfaiscau -> préfaisceau, misspelling

\textit{a morphism from $(U,T,H)$ to $(U',T',H')$ being a triple}
\[
(f_U, f_T, u) \qquad\qquad
\left\{\begin{array}{l}
\left.\begin{array}{llll}
f_U : & U & \longrightarrow & U' \\[1ex]
f_T : & T & \longrightarrow & T'
\end{array}\right\}\text{functors}\\[2ex]
u : H \longrightarrow H' \circ (f_U^\circ \times f_T),\\
\text{i.e. } H(u,t) \longrightarrow H'(f_U(u), f_T(t))\\
\text{a functorial homomorphism}\\
\text{(arrow in } (U \times T^\circ)^\wedge \simeq \uHom(U^\circ,T^\vee)\text{)),}
\end{array}\right.
\]
\textit{composition of arrows obvious.}

If $U \hookrightarrow X \hookleftarrow T$ is a triple, and
\[
f_{U'} : U' \longrightarrow U,\ f_{T'} : T' \longrightarrow T
\]
arrows in Cat, there is a privileged way to ``glue'' $U'$ and $T'$, in such a way that the object $(X',U',T')$ maps into $(X,U,T)$ inducing

\msp{page 61}

$f_{U'}$, $f_{T'}$:
\[
\begin{tikzcd}[column sep=large, row sep=large]
U' \arrow[r, hook, "i'"] \arrow[d, "f_{U'}"'] & X' \arrow[d, "f"] & T' \arrow[l, hook', "j'"'] \arrow[d, "f_{T'}"] \\
U \arrow[r, hook, "i"] & X & T, \arrow[l, hook', "j"']
\end{tikzcd}
\]
it is obtained with the help of the bifunctors
\[
\left\{\begin{aligned}
&H'(u',t') \overset{\text{def}}{=} H(f_{U'}(u'), f_{T'}(t'))\\
&\text{i.e. } H' = H \circ (f_{U'}^\circ \times f_{T'}).
\end{aligned}\right.
\]
This triple, i.e. this object of $\mathcal{T}$, represents the following functor on $\mathcal{T}$ :
\[
(X'',U'',T'') \longmapsto \mathrm{Hom}_{\mathrm{Cat}}(U'',U) \times \mathrm{Hom}_{\mathrm{Cat}}(T'',T),
\]
i.e. for every object $(X'',U'',T'')$ of $\mathcal{T}$ we have a canonical bijection, functorial in this object
\[
\mathrm{Hom}_{\mathcal{T}}((X'',U'',T''),(X',U',T')) \simeq \mathrm{Hom}_{\mathrm{Cat}}(U'',U) \times \mathrm{Hom}_{\mathrm{Cat}}(T'',T). \tag*{$(*)$}
\]

\printedpage{45}
% head: A. Grothendieck : Dérivateurs, Ch. XII, Édition des Universités de Montpellier II et Paris VII ; foot: Version 7 décembre 2002

No doubt we even have
\[
\uHom_{\mathcal{T}}((X'',U'',T''),(X',U',T')) \simeq \uHom_{\mathrm{Cat}}(U'',U) \times \uHom_{\mathrm{Cat}}(T'',T), \tag*{$(**)$}
\]
just as functorial, from which $(*)$ is deduced by passing to the sets of objects. We say that the triple $(X',U',T')$ is the triple \textit{induced} by the triple $(X,U,T)$, via the functors $f_{U'}$, $f_{T'}$.

\textbf{NB} If $f_{U'}$, $f_{T'}$ are fully faithful, then so is $f : X' \longrightarrow X$. We shall use this operation of privileged gluing mainly when $U' = U$, $f_{U'} = \mathrm{id}_U$, or conversely $T' = T$, $f_{T'} = \mathrm{id}_T$.

\subsubsection*{1.2\quad Relation with the gluing of topoi.}

The situation of a triple
\[
(\mathcal{X},\mathcal{U},\mathcal{T})
\]
of a topos $\mathcal{X}$, of an open $\mathcal{U}$ of $\mathcal{X}$, and of a

\msp{page 62}

complementary closed subtopos, is entirely similar (\textbf{NB} If $\mathcal{X} = \mathrm{Top}(X)$, with $X$ in Cat, then $\mathrm{Ouv}(\mathcal{X}) \simeq \mathrm{Ouv}(X)$, since the opens correspond to the subsheaves of the terminal sheaf on $\mathcal{X}$, i.e. of the terminal presheaf on $X$. This bijection is compatible with passage to complements.) Topos theory says that the pairs $(\mathcal{X},\mathcal{U},\mathcal{T})$ \texttt{[rather triples]} are reconstructed, up to an equivalence defined up to unique isomorphism, with the help of a triple
\[
\left\{\begin{aligned}
&\mathcal{U},\ \mathcal{T} \text{ two topoi}\\[1ex]
&\varphi : \mathrm{Fais}\,(\mathcal{U}) \longrightarrow \mathcal{T} \text{ a left exact functor (``accessible''),}
\end{aligned}\right.
\]
this functor $\varphi$ being obtained, when we start from
\[
\mathcal{U} \xrightarrow{\;i\;} \mathcal{X} \xleftarrow{\;j\;} \mathcal{T},
\]
as
\[
\varphi = j^*i_*.
\]
Moreover, the triple $(\mathcal{U},\mathcal{T},\varphi)$ can be chosen arbitrarily, with the sole restriction (practically harmless, since always satisfied for all the natural functors we encounter) that $\varphi$ be (besides being left exact) ``accessible''.

In the present case $\mathcal{X} = \mathrm{Top}(X)$, $\mathcal{U} = \mathrm{Top}(U)$, $\mathcal{T} = \mathrm{Top}(T)$, the functor $\varphi = j^*i_*$ is not only left exact, but it even commutes with arbitrary $\varprojlim$, since this is so (always) for $i_*$, and also (because the topos-theoretic $j$ is induced by the category-theoretic $j : T \longrightarrow X$) for $j^*$. In fact, we
\printedpage{46}
% head: A. Grothendieck : Dérivateurs, Ch. XII, Édition des Universités de Montpellier II et Paris VII ; foot: Version 7 décembre 2002
\msp{page 63}
find this :

\textbf{Proposition 2 :} \textit{Let us consider a situation of gluing two topoi $\mathcal{U}$, $\mathcal{T}$ by $\varphi : \mathrm{Fais}\,\mathcal{U} \longrightarrow \mathrm{Fais}\,\mathcal{T}$, a left exact functor. Let us suppose that $\mathcal{U} = \mathrm{Top}(U)$, $\mathcal{T} = \mathrm{Top}(T)$, with $U$ and $T$ in $\mathrm{Cat}$. For the glued topos $\mathcal{X}$ to be equivalent to a topos $\mathrm{Top}(X)$, with $X$ in $\mathrm{Cat}$, it is necessary and sufficient that the gluing functor $\varphi$ commute with arbitrary $\varprojlim$, or what amounts to the same thing, that it have a left adjoint :}
\[
\psi : \underbrace{\mathrm{Fais}\,\mathcal{T}}_{= T^\wedge} \longrightarrow \underbrace{\mathrm{Fais}\,\mathcal{U}}_{= U^\wedge}.
\]

We then expect that the gluing datum $\varphi$, or what amounts to the same thing the functor $\psi$ having a right adjoint, i.e. commuting with arbitrary $\varinjlim$, amounts to that of a gluing datum $H$
\[
H : U^\circ \times T \longrightarrow \mathrm{Ens}
\]
for $U$ and $T$. But we know well that the functors $\psi : T^\wedge \longrightarrow \mathcal{C}$ into a category $\mathcal{C}$ with arbitrary $\varinjlim$, which commute with the said $\varinjlim$, form a category equivalent (by the functor $\psi \longmapsto \psi|T(\subseteq T^\wedge)$) to that of functors
\[
\psi_0 : T \longrightarrow \mathcal{C} \quad\ (= \underbrace{U^\wedge}_{\uHom(U^\circ,\mathrm{Ens})} \text{ in this case}),
\]
\msp{page 64}
and giving such a functor $\psi_0$ does indeed amount to giving an $H$. Thus

\textbf{Corollary.} \textit{The category of topos-theoretic gluing data $\varphi : U^\wedge \longrightarrow T^\wedge$ between} $\mathrm{Top}(U)$, $\mathrm{Top}(T)$\textit{, such that the glued topos is a ``}paradisiacal\textit{ topos'', i.e. equivalent to a topos of the form $X^\wedge$ (with $X$ in} Cat\textit{), is equivalent to that of gluing data $H : U^\circ \times T \longrightarrow$} Ens \textit{between $U$ and $T$. The topos-theoretic datum $\varphi$ associated to $H$ is the functor right adjoint to $T^\wedge \longrightarrow U^\wedge$ deduced from $T \xrightarrow{H} U^\wedge$ by extension commuting with $\varinjlim$, or again it is given by}
\[
\varphi(\underbrace{F}_{\in U^\wedge})(t) \simeq \mathrm{Hom}_{U^\wedge}(H(t),F),
\]
\textit{where on the right-hand side $H$ is interpreted as a functor $T \longrightarrow U^\wedge$. The topos glued by $\varphi$ is canonically equivalent to} $\mathrm{Top}(X)$\textit{, where $X$ is the object of} Cat \textit{obtained by gluing with the help of $H$.}

From the point of view of sheaf theory, this gives the interpretation of $H$, as expressing the functor
\[
\varphi = j^*i_* : U^\wedge \longrightarrow T^\wedge
\]
by the preceding formula.

\printedpage{47}
% head: A. Grothendieck : Dérivateurs, Ch. XII, Édition des Universités de Montpellier II et Paris VII ; foot: Version 7 décembre 2002

\msp{page 65}

\subsubsection*{1.3\quad Arrows in $\mathrm{Cat}$, with source or target a glued object $(X,U,T)$.}

When $(X,U,T)$ is given by a gluing datum
\[
U,\ T,\ H : U^\circ \times T \longrightarrow \mathrm{Ens},
\]
the arrows
\[
f : X' \longrightarrow X
\]
in Cat are immediately described as follows : giving $f$ first defines
\[
U' = f^{-1}(U), \qquad\qquad T' = f^{-1}(T),
\]
a pair of an open of $X'$ and the complementary closed set, so that this datum is equivalent to that of

\begin{enumerate}[label={}]
\item[$1^\circ$)] a pair $(U',T')$ in $X'$
\item[$2^\circ$)] a morphism (in $\mathcal{T}$) from $(X',U',T')$ to $(X,U,T)$, which by 1.1 is equivalent to a triple
\[
\left\{\begin{array}{rccl}
f_{U'} : & U' & \longrightarrow & U\\[1ex]
f_{T'} : & T' & \longrightarrow & T\\[1ex]
u : & H_{U',T'} & \longrightarrow & H \circ (f_{U'}^\circ \times f_{T'}),
\end{array}\right.
\]
where we have put
\[
\left\{\begin{aligned}
&H_{U',T'} : U^{\prime\circ} \times T' \longrightarrow \mathrm{Ens}\\[1ex]
&H_{U',T'}(u',t') = \mathrm{Hom}_{X'}(u',t'),
\end{aligned}\right.
\]
i.e. $H_{U',T'}$ is the gluing datum which corresponds to the triple $(X',U',T')$ in $\mathcal{T}$. % [sic] conrrespond -> correspond, misspelling
\end{enumerate}

\msp{page 66}

We shall be more interested in the reverse situation, i.e. in the question of determining the functors
\[
\underbrace{X}_{= X_{U,T,H}} \quad \longrightarrow X'
\]
from $X$, obtained from $(U,T,H)$, to an arbitrary object of Cat. We find tautologically

\textbf{Proposition 3.} $\mathrm{Hom}_{\mathrm{Cat}}(X,X')$ \textit{is in canonical bijective correspondence, functorial in $X'$ in} Cat\textit{, with the set of triples}